\chapter{2009年上海卷（春）}

\section{填空题}

\begin{problem}

函数 \( y=\log_{2}(x-1) \) 的定义域是 \fillinblank{}.
\end{problem}

\begin{problem}

计算： \( (1-\i)^{2}= \) \fillinblank{}（i为虚数单位）.
\end{problem}

\begin{problem}

函数 \( y = \cos \frac{x}{2} \) 的最小正周期 T = \fillinblank{}.
\end{problem}

\begin{problem}

若集合 \( A = \{x \mid |x| > 1\} \)，集合 \( B = \{x \mid 0 < x < 2\} \)，则 \( A \cap B = \) \fillinblank{}. 
\end{problem}

\begin{problem}

抛物线 \( y^{2}=x \) 的准线方程是 \fillinblank{}.
\end{problem}

\begin{problem}

已知 \( \left|\bs{a}\right|=3,\quad\left|\bs{b}\right|=2 \) . 若 \( \bs{a}\cdot\bs{b}=-3 \) ，则 \( \bs{a} \) 与 \( \bs{b} \) 夹角的大小为 \fillinblank{}. 
\end{problem}

\begin{problem}

过点 \( A(4, -1) \) 和双曲线 \( \frac{x^{2}}{9} - \frac{y^{2}}{16} = 1 \) 右焦点的直线方程为 \fillinblank{}.
\end{problem}

\begin{problem}

在 \( \triangle ABC \) 中，若 \( AB = 3 \)， \( \angle ABC = 75^\circ \)， \( \angle ACB = 60^\circ \)，则 \( BC \) 等于 \fillinblank{}. 
\end{problem}

\begin{problem}

已知对于任意实数 \( x \)，函数 \( f(x) \) 满足 \( f(-x) = f(x) \). 若方程 \( f(x) = 0 \) 有 2009 个实数解，则这 2009 个实数解之和为 \fillinblank{}. 
\end{problem}

\begin{problem}

一只猴子随机敲击只有26个小写英文字母的练习键盘.若每敲1次在屏幕上出现一个字母，它连续敲击10次，屏幕上的10个字母依次排成一行，则出现单词“monkey”的概率为（结果用数值表示） \fillinblank{}.
\end{problem}

\begin{problem}

以下是面点师一个工作环节的数学模型：如图，在数轴上截取与闭区间 \( [0, 1] \) 对应的线段，对折后（坐标 1 所对应的点与原点重合）再均匀地拉成1个单位长度的线段，这一过程称为一次操作（例如在第一次操作完成后，原来的坐标 \( \frac{1}{4} \)、 \( \frac{3}{4} \)变成 \( \frac{1}{2} \)，原来的坐标 \( \frac{1}{2} \)变成1，等等）. 那么原闭区间 \( [0,1] \) 上（除两个端点外）的点，在第二次操作完成后，恰好被拉到与1重合的点所对应的坐标是 \fillinblank{}；原闭区间 \( [0,1] \) 上（除两个端点外）的点，在第 \( n \) 次操作完成后（\( n \geq 1 \)），恰好被拉到与1重合的点所对应的坐标为 \fillinblank{}.

\bitmapfigure[max width=0.3\linewidth,max totalheight=2.0cm]{img/2009/shanghai_spring/q11_fig1.png}

\end{problem}

\section{单选题}

\begin{problem}

在空间中，“两条直线没有公共点”是“这两条直线平行”的（\quad）

\choices
    {充分不必要条件.}
    {必要不充分条件.}
    {充要条件.}
    {既不充分也不必要条件.}

\end{problem}

\begin{answer}
B
\end{answer}
\begin{problem}

过点 \( P(0,1) \) 与圆 \( x^{2}+y^{2}-2x-3=0 \) 相交的所有直线中，被圆截得的弦最长时的直线方程是（\quad）

\choices
    {x=0.}
    {y=1.}
    {\( x + y - 1 = 0. \)}
    {\( x - y + 1 = 0. \)}

\end{problem}

\begin{answer}
C
\end{answer}
\begin{problem}

已知函数 \( f(x)=\begin{cases}3^{x+1}, & x \leq 0, \\ \log_{2} x, & x > 0.\end{cases} \) 若 \( f(x_{0})>3 \)，则 \( x_{0} \) 的取值范围是（\quad）

\choices
    {\( x_{0} > 8 \).}
    {\( x_{0} < 0 \) 或 \( x_{0} > 8 \).}
    {\( 0 < x_{0} < 8 \).}
    {\( x_{0} < 0 \) 或 \( 0 < x_{0} < 8 \).}

\end{problem}

\begin{answer}
A
\end{answer}
\begin{problem}

函数 \( y = 1 + \sqrt{1 - x^{2}} \)（\( -1 \leq x \leq 0 \)）的反函数图像是（\quad）

\bitmapfigure[max width=0.88\linewidth,max totalheight=6.2cm]{img/2009/shanghai_spring/q15_options.png}

\end{problem}

\begin{answer}
C
\end{answer}
\section{解答题}

\begin{problem}

如图，在斜三棱柱 \( ABC - A_1B_1C_1 \) 中， \( \angle A_1AC = \angle ACB = \frac{\pi}{2} \)， \( \angle AA_1C = \frac{\pi}{6} \)，侧棱 \( BB_1 \) 与底面所成的角为 \( \frac{\pi}{3} \)， \( AA_1 = 4\sqrt{3} \)， \( BC = 4 \). 求斜三棱柱 \( ABC - A_1B_1C_1 \) 的体积 \( V \). 

\bitmapfigure[max width=0.38\linewidth,max totalheight=4.8cm]{img/2009/shanghai_spring/q16_fig1.png}
\end{problem}

\begin{problem}

已知数列 \( \left\{a_{n}\right\} \) 的前 n 项和为 \( S_{n} \)， \( a_{1}=1 \)，且 \( 3a_{n+1}+2S_{n}=3 \)（ \( n \) 为正整数）. （1）求数列 \( \{a_{n}\} \)的通项公式； （2）记 \( S = a_1 + a_2 + \cdots + a_n + \cdots \). 若对任意正整数 \( n \)， \( kS \leq S_n \) 恒成立，求实数 \( k \) 的最大值. 
\end{problem}

\begin{problem}

我国计划发射火星探测器，该探测器的运行轨道是以火星（其半径 \( R = 34 \) 百公里）的中心 \( F \) 为一个焦点的椭圆. 如图，已知探测器的近火星点（轨道上离火星表面最近的点）\( A \) 到火星表面的距离为 \( 8 \) 百公里，远火星点（轨道上离火星表面最远的点）\( B \) 到火星表面的距离为 \( 800 \) 百公里. 假定探测器由近火星点 \( A \) 第一次逆时针运行到与轨道中心 \( O \) 的距离为 \( \sqrt{ab} \) 百公里时进行变轨，其中 \( a \)、\( b \) 分别为椭圆的长半轴、短半轴的长，求此时探测器与火星表面的距离（精确到 \( 1 \) 百公里）.

\bitmapfigure[max width=0.5\linewidth,max totalheight=4.8cm]{img/2009/shanghai_spring/q18_fig1.png}
\end{problem}

\begin{problem}

如图，在直角坐标系 xOy 中，有一组对角线长为 \( a_{n} \) 的正方形 \( A_{n}B_{n}C_{n}D_{n} \) （ \( n=1,2,\cdots \)）， 其对角线 \( B_n D_n \) 依次放置在 \( x \) 轴上（相邻顶点重合）. 设 \( \{a_n\} \) 是首项为 \( a \)，公差为 \( d(d > 0) \) 的等差数列，点 \( B_1 \) 的坐标为 \( (d, 0) \). (1) 当 a = 8, d = 4 时，证明：顶点 \( A_{1} \)、 \( A_{2} \)、 \( A_{3} \) 不在同一条直线上； （2）在（1）的条件下， 证明：所有顶点 \( A_{n} \) 均落在抛物线 \( y^{2}=2x \) 上； （3）为使所有顶点 \( A_{n} \) 均落在抛物线 \( y^{2}=2px \) （p>0）上， 求 a 与 d 之间所应满足的关系式.

\bitmapfigure[max width=0.54\linewidth,max totalheight=5.0cm]{img/2009/shanghai_spring/q19_fig1.png}
\end{problem}

\begin{problem}

设函数 \( f_{n}(\theta)=\sin^{n}\theta+(-1)^{n}\cos^{n}\theta,\quad0\leq\theta\leq\frac{\pi}{4} \)，其中 n 为正整数. （1）判断函数 \( f_{1}(\theta) \)、 \( f_{3}(\theta) \) 的单调性，并就 \( f_{1}(\theta) \) 的情形证明你的结论； （2）证明： \( 2f_{6}(\theta)-f_{4}(\theta)=\left(\cos^{4}\theta-\sin^{4}\theta\right)\left(\cos^{2}\theta-\sin^{2}\theta\right) \); （3）对于任意给定的正整数 n，求函数 \( f_{n}(\theta) \) 的最大值和最小值.
\end{problem}

